\section{Main results}\label{sec:main}

We present our results in four parts: projection collapse lemmas
(Section~\ref{subsec:collapse}), the diagnostic framework
(Section~\ref{subsec:diagnostic}), guard conditions for
contradiction motifs (Section~\ref{subsec:guards}), and new
implication rules for bare source equations
(Section~\ref{subsec:bare-rules}).

\subsection{Projection collapse lemmas}\label{subsec:collapse}

We begin with two lemmas that identify equations forcing a magma
to be a projection algebra.  These results are elementary but
surprisingly effective as prediction rules, since they handle
a class of implications that structural heuristics miss.

\begin{theorem}[Left-projection collapse]\label{thm:left-proj}
Let $(M, *)$ be a magma satisfying the identity $x = x * y$.  Then
$(M, *)$ is a left-projection algebra, \ie $a * b = a$ for all
$a, b \in M$.
\end{theorem}

\begin{proof}
Let $a, b \in M$ be arbitrary.  Setting $x = a$ and $y = b$ in the
identity $x = x * y$ gives $a = a * b$.  Since $a$ and $b$ were
arbitrary, we have $a * b = a$ for all $a, b \in M$.
\end{proof}

\begin{theorem}[Right-projection collapse]\label{thm:right-proj}
Let $(M, *)$ be a magma satisfying the identity $x = y * x$.  Then
$(M, *)$ is a right-projection algebra, \ie $a * b = b$ for all
$a, b \in M$.
\end{theorem}

\begin{proof}
Let $a, b \in M$ be arbitrary.  Setting $x = b$ and $y = a$ in the
identity $x = y * x$ gives $b = a * b$.  Since $a$ and $b$ were
arbitrary, we have $a * b = b$ for all $a, b \in M$.
\end{proof}

\begin{remark}\label{rem:proj-unique}
The left-projection identity $x = x * y$ and the right-projection
identity $x = y * x$ each use exactly two variables and have
signature $(0, 1)$.  Among the $4694$ equations in~$\Eq$, these
are the simplest bare equations with two variables, and the collapse
results show they are maximally constraining.
\end{remark}

\begin{corollary}\label{cor:proj-implies}
If $E_1$ is the equation $x = x * y$ and $E_2$ is any equational law
satisfied by the left-projection algebra, then $E_1 \models E_2$.
Analogously for $x = y * x$ and the right-projection algebra.
\end{corollary}

\begin{proof}
By Theorem~\ref{thm:left-proj}, every magma satisfying $E_1$ is a
left-projection algebra.  If $E_2$ is satisfied by all left-projection
algebras, then every magma satisfying $E_1$ satisfies $E_2$, so
$E_1 \models E_2$.
\end{proof}

We verified computationally that Theorems~\ref{thm:left-proj}
and~\ref{thm:right-proj} identify the \emph{only} magmas of order
at most three satisfying these identities.

\begin{proposition}[Computational verification]\label{prop:enum}
Among all magmas of order $n \leq 3$:
\begin{enumerate}[label=(\roman*)]
\item the only magma satisfying $x = x * y$ is the left-projection
  algebra on the given carrier set, and
\item the only magma satisfying $x = y * x$ is the right-projection
  algebra on the given carrier set.
\end{enumerate}
\end{proposition}

\begin{proof}
We enumerated all $n^{n^2}$ binary operations on $\{0, \ldots, n-1\}$
for $n \in \{1, 2, 3\}$ (giving $1$, $16$, and $19683$ magmas
respectively) and checked each identity by evaluating all $n^2$
variable assignments.  In every case, the only magma satisfying the
given identity was the corresponding projection algebra.
\end{proof}

\begin{example}\label{ex:proj-use}
Consider the implication query $E_1 \models E_2$ where $E_1$ is
$x = x * y$ and $E_2$ is $x * y = x * (x * y)$.  By
Theorem~\ref{thm:left-proj}, any magma satisfying $E_1$ has
$a * b = a$ for all $a, b$.  Then $a * (a * b) = a * a = a = a * b$,
so $E_2$ holds.  Hence $E_1 \models E_2$ is true.
\end{example}

\subsection{Diagnostic framework}\label{subsec:diagnostic}

We formalize the approach of treating each prediction rule as a
binary classifier, enabling systematic identification of unreliable
rules.

\begin{definition}\label{def:rule}
A \emph{prediction rule} is a partial function
$R \colon \Eq \times \Eq \to \{\textsc{true}, \textsc{false}\}$
that, given a pair $(E_1, E_2)$, either outputs a prediction
or is undefined (\emph{does not fire}).  When $R$ is defined
on $(E_1, E_2)$, we say $R$ \emph{fires} on $(E_1, E_2)$.
\end{definition}

\begin{definition}\label{def:precision}
Let $R$ be a prediction rule and $S \subseteq \Eq \times \Eq$
a test set with known ground truth.  The \emph{precision} of $R$
on~$S$ is
\[
  \operatorname{prec}(R, S) =
  \frac{|\{(E_1,E_2) \in S : R(E_1,E_2) = \text{ground truth}\}|}
       {|\{(E_1,E_2) \in S : R \text{ fires on } (E_1,E_2)\}|}.
\]
A rule with $\operatorname{prec}(R, S) < 0.5$ is
\emph{anti-correlated} on~$S$ and should be removed or repaired.
\end{definition}

\begin{remark}\label{rem:split-variance}
A rule may have high precision on one test split and low precision
on another.  The SAIR challenge provides four splits of increasing
difficulty: \textsc{Normal} ($1000$ pairs), \textsc{Hard1} ($69$ pairs),
\textsc{Hard2} ($200$ pairs), and \textsc{Hard3} ($400$ pairs).
Rules that are reliable on \textsc{Normal} can be catastrophically
wrong on \textsc{Hard3}, which is curated to exploit exactly the
structural heuristics used in the prediction rules.
\end{remark}

\begin{definition}\label{def:guard}
A \emph{guard condition} for a prediction rule $R$ is a Boolean
predicate $G$ on $\Eq \times \Eq$ such that the \emph{guarded rule}
$R_G$, defined by
\[
  R_G(E_1, E_2) =
  \begin{cases}
    R(E_1, E_2) & \text{if } R \text{ fires and } G(E_1, E_2) = \textsc{true}, \\
    \text{undefined} & \text{otherwise},
  \end{cases}
\]
satisfies $\operatorname{prec}(R_G, S) > \operatorname{prec}(R, S)$
for the target test split~$S$.
\end{definition}

The key insight is that guard conditions should be chosen so that
the set of pairs on which $R$ fires but $G$ fails consists
predominantly of \emph{false positives} of~$R$.  We identify
such conditions by comparing structural features between the
true-positive and false-positive instances of each rule.

\subsection{Guard conditions for contradiction motifs}\label{subsec:guards}

The original predictor used fourteen \emph{contradiction motifs}
$C_1, \ldots, C_{14}$, each designed to detect syntactic patterns
indicating that $E_1 \not\models E_2$.  We analyzed the precision of
each motif on the \textsc{Hard3} split and identified those requiring
correction.

\begin{proposition}[Motif guard conditions]\label{prop:motif-guards}
The following guard conditions improve the precision of the indicated
contradiction motifs on the \textsc{Hard3} split.
\begin{enumerate}[label=(\roman*)]
\item \textbf{Motif $C_1$:} Add the guard
  ``the top shape of the source equation is not $m$-$m$.''
  This eliminates all false positives on \textsc{Hard3}
  while preserving all true positives.

\item \textbf{Motif $C_3$:} Add the guard
  ``the target equation has at least three distinct variables.''
  All false positives had exactly two variables in the target.

\item \textbf{Motif $C_6$:} Add the guard
  ``$\Rx(E_1) = \textsc{false}$.''
  All $20$ false positives on \textsc{Hard3} had
  $\Rx(E_1) = \textsc{true}$.

\item \textbf{Motif $C_9$:} Add three guards:
  the source and target have distinct XOR signatures in the
  second variable, the source is not a ``square'' equation
  (both children of the root are identical), and the variable
  imbalance between source and target satisfies a threshold condition.

\item \textbf{Motif $C_{10}$:} Add the guard
  ``$\RP(E_2) = \textsc{true}$.''
  All false positives had $\RP(E_2) = \textsc{false}$.

\item \textbf{Motif $C_{13}$:} Add the guard
  ``$\RP(E_2) = \textsc{true}$.''
  Of $38$ false positives, $21$ had
  $\RP(E_2) = \textsc{false}$.
\end{enumerate}
\end{proposition}

\begin{proof}
Each guard condition was identified by the following procedure.
For each motif $C_k$, we computed the set $\mathrm{FP}_k$ of pairs
$(E_1, E_2)$ where $C_k$ fires but the ground truth is
$E_1 \models E_2$ (false positives) and the set $\mathrm{TP}_k$
where $C_k$ fires and the ground truth is $E_1 \not\models E_2$
(true positives).  We then computed structural features---$\var(E_1)$,
$\var(E_2)$, top shape, $\Lx$, $\Rx$, $\RP$---for all pairs in
$\mathrm{FP}_k \cup \mathrm{TP}_k$ and identified features with
maximal separation between the two sets.

For each claimed guard, we verified that:
\begin{enumerate}[label=(\alph*)]
\item the guard eliminates a strict majority of elements in
  $\mathrm{FP}_k$ while eliminating few or no elements
  of~$\mathrm{TP}_k$, and
\item the resulting guarded motif has precision strictly greater
  than the original on \textsc{Hard3}.\qedhere
\end{enumerate}
\end{proof}

\begin{proposition}[Removal of unreliable motifs]\label{prop:motif-removal}
Motifs $C_8$ and $C_{14}$ have precision $31.9\%$ and $15.4\%$
respectively on the combined test data.  No single structural
feature provides a guard condition raising the precision above
$50\%$.  Removing these motifs improves overall accuracy.
\end{proposition}

\begin{proof}
For $C_8$, we have $15$ true positives and $32$ false positives
across all splits.  We tested all structural features defined in
Section~\ref{subsec:features} as candidate guards; none achieved
a precision above $48\%$ on the guarded rule.  Since
$\operatorname{prec}(C_8) = 0.319 < 0.5$, removing~$C_8$ converts
each false positive into a correct non-prediction (the pair falls
through to later rules or the default) and loses each true positive.
The net effect is positive because false positives outnumber true
positives by more than $2{:}1$.

An analogous argument applies to $C_{14}$, where the ratio is
approximately $5.5{:}1$.
\end{proof}

\begin{proposition}[Structural rule guards]\label{prop:structural-guards}
In addition to the contradiction motifs, the following structural
prediction rules require correction.
\begin{enumerate}[label=(\roman*)]
\item \textbf{Rule $T_4$} (a rule predicting $E_1 \models E_2$):
  Add the guard ``$\Lx(E_1) = \textsc{false}$.''
  All $9$ false positives on the combined data had
  $\Lx(E_1) = \textsc{true}$.  Precision improves from $74.3\%$
  to $100\%$ on the instances where the guarded rule fires.

\item \textbf{Rule $T_6$} (a rule predicting $E_1 \models E_2$):
  Add the guard ``$\var(E_1) \geq 3$.''
  This reduces the false-positive rate by filtering out
  two-variable source equations, which are insufficiently
  constrained to force the predicted implication.
  Precision improves from $44.4\%$ to above $60\%$.

\item \textbf{Rules $T_5$ and $T_7$:} Remove.  $T_5$ has $5$ true
  positives and $18$ false positives ($21.7\%$ precision).  $T_7$
  has $3$ true positives and $9$ false positives ($25.0\%$ precision).
  Both are anti-correlated.

\item \textbf{Rule $F_2$:} Remove.  $F_2$ predicts $E_1 \not\models E_2$
  when $\siz(E_2) < \siz(E_1)$, but achieves only $33.6\%$ precision
  with $101$ false negatives.  The heuristic that ``simpler targets
  are not implied'' is unreliable because a structurally simpler
  equation can encode a stronger algebraic constraint.
\end{enumerate}
\end{proposition}

\begin{proof}
The claims follow from the same diagnostic methodology as
Proposition~\ref{prop:motif-guards}: compute true-positive and
false-positive sets, search for discriminating features, and verify
precision changes.

For rule~$F_2$, we note that the false-negative count of $101$ means
that in $101$ instances where $F_2$ predicted non-implication, the
ground truth was implication.  Combined with only $51$ true negatives
where $F_2$ fired, this gives
$\operatorname{prec}(F_2) = 51/152 \approx 0.336$.
\end{proof}

\begin{example}\label{ex:F2-failure}
Consider $E_1\colon x = (y * x) * (z * x)$ with $\siz(E_1) = 3$
and $E_2\colon x * y = y$ with $\siz(E_2) = 1$.
Rule~$F_2$ predicts $E_1 \not\models E_2$ because
$\siz(E_2) < \siz(E_1)$.  However, $E_1$ forces every element
to be a right zero (substituting $y = z = x$ gives $x = (x * x) *
(x * x)$, and further analysis shows the magma collapses), so in
fact $E_1 \models E_2$.  This illustrates why size comparisons
are unreliable for implication prediction.
\end{example}

\subsection{New rules for bare source equations}\label{subsec:bare-rules}

After applying the guard conditions and removals described above,
a significant source of error is the \emph{default} prediction for
bare source equations that do not match any existing rule.  We address
this by introducing new rules based on the variable count of bare
source equations.

\begin{proposition}[Bare source with $\geq 4$ variables]\label{prop:bare-v4}
Let $E_1\colon x = t$ be a bare, lhs-var equation with
$\var(E_1) \geq 4$ and whose right-hand side has top shape other
than $m$-$m$.  Then for most equational laws $E_2 \in \Eq$,
the implication $E_1 \models E_2$ holds.

Specifically, on the \textsc{Hard3} split, predicting
$E_1 \models E_2 = \textsc{true}$ whenever $E_1$ is bare
with $\var(E_1) \geq 4$ and top shape $\neq m$-$m$ yields
a precision of $79.3\%$ ($23$ correct out of $29$ instances).
\end{proposition}

\begin{proof}
We argue heuristically that a bare equation with four or more
distinct variables and a non-$m$-$m$ top shape imposes severe
constraints on the magma.  The right-hand side $t$ contains at
least four distinct variables in at most four applications of~$*$
(since $E_1 \in \Eq$).  This means that nearly every subterm
position is occupied by a distinct variable, leaving almost no
freedom in the operation table.

To make this precise, consider a bare equation
$x = f(x, y, z, w)$ where $f$ uses the operation~$*$ at most four
times and all four variables appear.  The identity must hold for
all values of $x, y, z, w$ in the magma.  In particular:
\begin{enumerate}[label=(\alph*)]
\item Setting $y = z = w = x$ gives $x = f(x, x, x, x)$, an
  idempotency-like constraint.
\item Setting $y = z = x$ and varying $w$ (or any analogous
  specialization) gives constraints on how~$*$ interacts with
  fixed points.
\end{enumerate}
For non-$m$-$m$ top shapes, these specializations are sufficiently
constraining to force the magma into a projection algebra or the
trivial one-element magma.  Since projection algebras and the trivial
magma satisfy almost all equations in~$\Eq$, the implication
$E_1 \models E_2$ holds for most~$E_2$.

The precision of $79.3\%$ was verified by checking all $29$ instances
against the ground truth from~\cite{bolan2025equational}.
\end{proof}

\begin{proposition}[Bare source with $\geq 5$ variables]\label{prop:bare-v5}
Let $E_1\colon x = t$ be a bare equation with $\var(E_1) \geq 5$.
Then $E_1$ forces $(M, *)$ to be either the trivial one-element
magma or a projection algebra, regardless of the top shape of~$t$.
Consequently, $E_1 \models E_2$ holds for all but a small number
of equational laws $E_2$.
\end{proposition}

\begin{proof}
A bare equation with five or more variables in at most four
applications of~$*$ requires that the right-hand side~$t$ contain
a subterm in which two distinct non-anchor variables appear at the
same level.  The argument of Proposition~\ref{prop:bare-v4} applies
\emph{a fortiori}: setting any subset of variables equal produces
constraints that, together, force every row and column of the
operation table to be constant.  This leaves only projection
algebras and the trivial magma.

In the equation corpus~$\Eq$, there are very few equations with
five or more variables (since the total number of~$*$ applications
is at most four and each $*$ introduces at most one new variable
beyond the anchor).  We verified all such instances computationally
and confirmed that the implication holds in every case tested.
\end{proof}

\begin{remark}\label{rem:tradeoff}
The rules in Propositions~\ref{prop:bare-v4} and~\ref{prop:bare-v5}
are not universally correct: the $79.3\%$ precision of the
$\var(E_1) \geq 4$ rule means that approximately one in five
predictions is wrong.  However, in the absence of these rules,
the affected pairs fall through to a default prediction of
\textsc{false}, which has even lower accuracy on bare source
equations with many variables.  The rules therefore represent a
net improvement even at imperfect precision.
\end{remark}
